% ==============================================================================
% Section: Policy Alternatives
% Status: COMPLETE (cut to ~1.5 pages per Revision_Instructions_2, 2026-03-28)
% ==============================================================================

\section{Policy Alternatives}
\label{sec:policy_alternatives}

The analysis in Sections~\ref{sec:allocative_cost}--\ref{sec:21billion}
shows that the Shapiro settlement reduces equilibrium capacity prices by
\$4--\$8/MW-day at the RTO level without altering the underlying markup
structure ($L = 0.54$ under either the settlement or the unconstrained
equilibrium). This section considers four structural alternatives that
address the conditions supporting cap-binding outcomes (concentrated
supply and limited competitive fringe) rather than suppressing the price
symptom.

\textbf{Option~1: The Shapiro Settlement (current policy).}
The settlement reshapes the VRR demand curve for 2026/27 and 2027/28,
providing immediate price relief without structural changes to the seller
side. Its key advantages are speed and regulatory certainty under FPA~\S205.
Its key limitation is that it does not reduce the number of strategic sellers,
expand the competitive fringe, or steepen the VRR demand curve, leaving the
$K \leq 3$ regime and the associated markup structure intact.

\textbf{Option~2: Transmission expansion (CETL).}
Increasing fringe supply $\bar{Q}_f$ reduces equilibrium prices organically
by shifting the market-clearing condition without modifying the demand curve.
A 30\% increase in competitive fringe supply at $K = 3$ reduces the
price-cost margin from 0.54 to approximately 0.30 in 2026/27, a 24-percentage-point
reduction that exceeds the settlement's zero-point reduction in markup.
The LDA cross-section (Section~\ref{subsec:res_lda}) confirms: a
10-percentage-point increase in import penetration is associated with a
4--6-point reduction in the local price-cost margin. The relevant policy
instrument is PJM's Regional Transmission Expansion Plan directed toward
import-constrained LDAs.

\textbf{Option~3: VRR curve redesign.}
Steepening the VRR demand curve reduces equilibrium markups through the
ODE mechanism: larger $|D'(p)|$ forces a more competitive supply function
at every price below the cap. A slope scaling factor of approximately 1.4
or greater would shift the 2026/27 outcome from cap-binding to
interior-clearing at $K = 3$, allowing prices to operate freely below
Point~(a). FERC Order ER25-1357 moved in this direction by steepening
the sloped segment. Further recalibration is available under FERC's
ongoing VRR review.

\textbf{Option~4: Entry promotion and deconcentration.}
The transition from cap-binding to interior equilibria occurs between
$K = 3$ and $K = 4$: adding a fourth symmetric strategic seller reduces
the price-cost margin by 20--33 percentage points and moves the 2026/27
equilibrium from \$329 to \$228 (Section~\ref{subsec:res_concentration}).
The relevant levers are interconnection queue reform (FERC Orders 2023 and
2023-A), expanded demand response and battery storage participation, and
targeted offer price mitigation or divestiture for the most capacity-concentrated
LDAs (PS, PS~NORTH, COMED, EMAAC).

\begin{table}[ht]
  \centering
  \small
  \caption{Comparative Assessment of Policy Alternatives}
  \label{tab:policy_comparison}
  \begin{tabular}{p{3.2cm}p{2.8cm}lp{2cm}p{3cm}}
    \hline\hline
    Option & Price reduction mechanism & Duration & Signal preserved & Regulatory lever \\
    \hline
    1. Settlement (current) & Cap/floor reshape VRR & 2 years & No (cap binding) & FERC \S205 tariff \\
    2. CETL expansion       & Competitive fringe $\uparrow$ & Permanent & Yes & RTEP, FERC Tx planning \\
    3. VRR redesign         & Demand elasticity $\uparrow$ & Permanent & Yes & FERC VRR recalibration \\
    4. Entry/deconcentration & $K \uparrow$ / fringe $\uparrow$ & Permanent & Yes & FERC Ords. 2023/2023-A; mitigation \\
    \hline\hline
  \end{tabular}
  \smallskip

  \noindent\footnotesize
  \textit{Notes:} Signal preservation indicates whether the capacity clearing
  price reflects unconstrained strategic equilibrium (Options~2--4) or is
  truncated by an administrative cap (Option~1). Duration refers to whether
  the price reduction persists beyond a fixed delivery window.
\end{table}

Under the model's assumptions, Options~2--4 dominate Option~1 along every
structural dimension: they produce permanent price reductions by addressing
the structural conditions associated with cap-binding outcomes rather than
the symptom. The settlement is justified as an emergency measure during a
period of acute schedule disruption, but it should not substitute for the
structural reforms that would render a similar intervention unnecessary in
future auctions. PJM's own settlement framing calls for ``reform the
capacity market'' alongside the cap \citep{PJM2025_settlement_slides},
consistent with this ranking.
